\subsection{\problemblock{*Optimize} Data Block}\label{sec:block-optimize}

Parametric optimization can be used to maximize or minimize a trajectory
characteristic, such as final range or velocity, and it can be used to satisfy
trajectory constraints, such as a final impact position or flight-path angle.
Although optimization is one of the most powerful features of TAOS, the numerical
optimization methods are complex and are not as reliable as the ordinary trajectory
integration methods.

A parametric optimization problem varies a set of input parameters
$\boldsymbol{x}$ to maximize or minimize an objective function
$f(\boldsymbol{x})$ subject to a set of constraints
$\boldsymbol{c}(\boldsymbol{x})$.  The parameters are numerical values in the
problem file.  The objective is a trajectory output variable evaluated at a specified
point, usually the end of a segment.  Constraints may be equalities, for which an
output variable must equal a value, or inequalities, for which a limit or boundary is
imposed.

TAOS uses the nonlinear constrained recursive-quadratic-programming method
\texttt{vf02ad}, described in \cref{sec:optimization}.  The partial
derivatives required by the algorithm are estimated with forward or central
differences.

\taossubhead{Optimization Loops}

Optimization is often applied to only part of a trajectory.  This part is called the
\emph{optimization loop}.  The loop begins with the first segment containing one of
the optimization parameters and extends to the segment at which the objective
function is evaluated.  It may contain any number of segments.

A problem may contain as many as five optimization loops, identified by the letters
\texttt{a} through \texttt{e}.  The loops may not overlap and may not be nested.
Each must be completely independent of every other search or optimization loop.
The loop letter associates the parameter placeholders with the
\problemblock{*optimize} block that defines the objective and constraints.

An optimization loop may have any number of parameters, but in practice the method
works best with fewer than about 25.  It normally converges faster and more reliably
when fewer parameters are used, so the objective is to include only enough
parameters to represent the trajectory adequately.  Parameter numbers start at one
and must be sequential, although the placeholders need not appear in numerical order
in the problem file.

\taossubhead{Optimization Parameters}

An input value is designated as parameter $n$ of loop $x$ by replacing the number
with \texttt{optx-n}.  For example,

\begin{lstlisting}[style=taosmanual]
*fly alpha=opta-2
\end{lstlisting}

makes angle of attack parameter 2 in loop \texttt{a}.  Any input variable that would
otherwise contain a number, rather than a table name or keyword, can be used as an
optimization parameter.  Initial conditions and guidance rules are the most common
choices.

A frequently used technique for optimizing the shape of a trajectory is to define a
segment with a tabulated guidance rule and an optimized final condition:

\begin{lstlisting}[style=taosmanual]
*fly alpha vrs tseg
  opta-1   0.0
  opta-2   5.0
  opta-3  10.0
  opta-4  15.0
  opta-5  opta-6

*when tseg>opta-6 stop
\end{lstlisting}

Six parameters define this angle-of-attack time history.  The segment times must be
estimated, but optimization determines the angle-of-attack values required to satisfy
the objective and constraints.  The final time is also made a parameter to provide
additional flexibility in shaping the flight path.

\taossubhead{Objective Function}

The optimization parameters alone do not define the loop.  The objective and any
constraints are entered in the \problemblock{*optimize} block.  The first line of the
block defines the objective function.  Time, range, and velocity are commonly used,
but any output variable is valid.  For example,

\begin{lstlisting}[style=taosmanual]
*optimize a for vel=max on segment 6, trajectory 1
\end{lstlisting}

sets the objective to maximum velocity at the end of segment 6 in trajectory 1.  The
loop letter follows the \problemblock{*optimize} keyword, and the keyword
\texttt{for} is required.  The output variable is set equal to \texttt{min} or
\texttt{max}.  The segment and trajectory numbers are required.

\taossubhead{Constraints}

The next lines may define equality or inequality constraints.  Constraints are
optional.  For example,

\begin{lstlisting}[style=taosmanual]
constrain vel=3000
constrain gamgd=-80
constrain qmin<0
\end{lstlisting}

constrains trajectory output variables, not the input parameters.  Because these lines
do not specify endpoints, they are evaluated at the same segment and trajectory as
the objective.  In continuation of the preceding example, velocity must equal 3000,
geodetic flight-path angle must equal $-80^\circ$, and the user-defined variable
\texttt{qmin} must remain below zero at the end of segment 6 in trajectory 1.

A constraint may name a different endpoint:

\begin{lstlisting}[style=taosmanual]
constrain vel on segment 3, trajectory 2 = 3000
\end{lstlisting}

which constrains velocity at the end of segment 3 in trajectory 2 to 3000 ft/s.

Integral user-defined variables are often used for inequality constraints that must
hold over the entire flight path.  For example, a vehicle with aerodynamic controls
may require a minimum dynamic pressure of
500~lbf/ft\textsuperscript{2}.  A suitable integral variable is

\begin{align}
  q_{\min} &= \int (q-500)^2\,dt,
    && q<500, \label{eq:optimization-qmin-violation}\\
  q_{\min} &= 0,
    && \text{otherwise}. \label{eq:optimization-qmin-satisfied}
\end{align}

Constraining \texttt{qmin<0} forces the trajectory to remain above the minimum
pressure.  Altitude, load factors, and other flight-path limits can be handled in the
same way.  Integral variables are defined as described in
\cref{sec:trajectory-block-define}.

\taossubhead{Multiple-Vehicle Constraints}

A complete constraint consists of the keyword \texttt{constrain}, a relationship,
and an optional reference factor.  The relationship contains an output-variable name,
one of the relations \texttt{=}, \texttt{<}, or \texttt{>}, and a value.  The value
may be a number or another output variable evaluated at the same or a different
trajectory point.  This feature is useful for intercept calculations.  For example,

\begin{lstlisting}[style=taosmanual]
constrain alt=alt[2]
\end{lstlisting}

requires the altitude on the optimization trajectory to equal the altitude on trajectory
2.  A bracketed subscript is one way to enter a trajectory number.

The same type of constraint can be written with keywords:

\begin{lstlisting}[style=taosmanual]
constrain alt on segment 6, trajectory 1
  = alt on segment 4, trajectory 3
\end{lstlisting}

This requires the altitude at the end of segment 6 in trajectory 1 to equal the altitude
at the end of segment 4 in trajectory 3.  Both endpoints must have been computed by
the time the objective is evaluated.  The keywords \texttt{on segment} and
\texttt{trajectory} supply the endpoint numbers.

Trajectory subscripts and endpoint keywords are alternative syntaxes:

\begin{lstlisting}[style=taosmanual]
constrain east[2] on segment 4 = east[1] on segment 7

constrain east on segment 4, trajectory 2
  = east on segment 7, trajectory 1
\end{lstlisting}

Both forms define the same relationship.  A trajectory number may not be supplied
both as a subscript and with the \texttt{trajectory} keyword.  If both segment and
trajectory numbers are given, the segment number comes first.  If neither is given,
the objective-function endpoint is used.

\taossubhead{Referencing Values}

Optimization is more reliable when the parameters, objective, and constraints have
approximately the same order of magnitude.  If some values are orders of magnitude
larger than others, numerical difficulties often result.  TAOS therefore provides
reference factors that normalize values during the calculation.

A constraint containing a nonzero number, such as \texttt{alt=30000}, is
automatically referenced by that number.  Other constraints may supply an explicit
factor:

\begin{lstlisting}[style=taosmanual]
constrain vel on segment 11 = vel[3] on segment 4, ref=1000
\end{lstlisting}

Both velocities are then divided by 1000 for the numerical calculation.  The
\texttt{ref} item is optional, must be the last item in the constraint, and does not
change the physical relationship.

\taossubhead{Optimization Control Variables}

After the constraints, optional control variables may be entered.  They are listed in
\cref{tab:optimization-control-variables}.  All have reasonable defaults.
The objective reference \statevar{fref} normalizes the objective in the same way that
\texttt{ref} normalizes a constraint.  For best results,
$|f(\boldsymbol{x})/f_{\mathrm{ref}}|$ should be between 1 and 10.  The convergence
accuracy \statevar{tol} determines when the method considers a solution valid.
Smaller values can produce a better solution, but require more iterations; if the
value is too small, convergence may not be achieved.

\begin{table}[htbp]
  \centering
  \caption{Optimization data-block input variables.}
  \label{tab:optimization-control-variables}
  \small
  \begin{tabularx}{0.98\linewidth}{@{}>{\ttfamily}lX>{\ttfamily}lX@{}}
    \toprule
    \normalfont Variable & Description & \normalfont Variable & Description \\
    \midrule
    fref & Objective-function reference value (default 1.0) &
    derivs & Derivative method (default 1) \\
    tol & Convergence accuracy (default $1.0\times10^{-6}$) &
    dx & Derivative increment (default $1.0\times10^{-8}$) \\
    maxitr & Maximum number of iterations (default 50) &
    adjust & Guidance-table adjustment (default 0) \\
    integ & Integration control (default 1) &
    surveys & Survey-control flag (default 0) \\
    restarts & Restart control (default 0) &
    print & Printout control (default 0) \\
    \bottomrule
  \end{tabularx}
\end{table}

The maximum iteration count \statevar{maxitr} controls the number of calls to
\texttt{vf02ad}.  It is related to, but is not the same as, the number of objective
function evaluations or trajectories computed.  With \texttt{maxitr=0}, one
trajectory is calculated using the initial parameter estimates.  This is useful for
checking whether the starting trajectory is reasonable.

Optimization recomputes trajectories many times.  The integration control
\statevar{integ} determines whether all trajectories in the problem are recomputed
for every function evaluation (1) or only the trajectory containing the objective is
recomputed (0).  If the problem is restricted to one trajectory, zero can reduce the
run time substantially.  If constraints refer to other trajectories, one is normally
necessary.

\taossubhead{Automatically Restarting Optimization}

If the method has not converged within the maximum number of iterations, it can be
restarted from the current solution.  A restart sometimes permits further progress
toward the minimum.  The variable \statevar{restarts} specifies how many restarts
are allowed; zero means none.  When restarts are used, \statevar{maxitr} should be
reduced to a value such as 20.  On the last restart, TAOS allows twice the stated
number of iterations in an effort to achieve convergence.

When an angle-of-attack time history is optimized, the segment times must be
estimated.  The method works best when the points are evenly distributed:

\begin{lstlisting}[style=taosmanual]
*fly alpha vrs tseg
  opta-1    0
  opta-2   20
  opta-3   40
  opta-4   60
  opta-5   80
  opta-6  100
  opta-7  opta-8

*when tseg>opta-8 stop
\end{lstlisting}

During optimization the final segment time \texttt{opta-8} can become either very
large or very close to the preceding point at 100 seconds, producing an uneven
distribution that affects the solution.  The guidance-table adjustment variable
\statevar{adjust} controls whether the fixed time points are redistributed evenly at
each restart.  Zero disables the adjustment; one enables it.

\taossubhead{Numerical Derivatives}

The optimization algorithm requires partial derivatives.  They are computed by a
central-difference method when \texttt{derivs=0}, or a forward-difference method
when \texttt{derivs=1}.  Forward differences use

\begin{equation}
  \frac{\partial f}{\partial x}
  \approx \frac{f(x+\delta x)-f(x)}{\delta x},
  \label{eq:optimization-forward-difference}
\end{equation}

and central differences use

\begin{equation}
  \frac{\partial f}{\partial x}
  \approx \frac{f(x+\delta x)-f(x-\delta x)}{2\delta x}.
  \label{eq:optimization-central-difference}
\end{equation}

The derivative increment \statevar{dx} controls $\delta x$.  Smaller increments
improve the derivative until numerical precision is reached.  Central differences
require more function evaluations, but usually give more accurate derivatives than
forward differences for the same increment.  The defaults normally work well and
do not need to be changed.

\taossubhead{Survey Control Flag}

If surveys are used with optimization, \statevar{surveys} controls the starting
trajectory for each set of survey values.  The first survey case always starts from the
\texttt{par-n} values in the \problemblock{*optimize} block.  With the default
\texttt{surveys=0}, the same starting trajectory is used for every later survey case.
With \texttt{surveys=1}, the optimum from the preceding case becomes the starting
trajectory for the next.  This often improves convergence when the optimum changes
only slightly between adjacent survey values.

\taossubhead{Printout Control}

TAOS always prints an optimization summary giving the parameters, objective, and
constraints at each iteration.  If \texttt{print=1}, the complete trajectory is also
printed after every iteration.  This can help locate a problem, but it produces a large
amount of output and is not recommended unless \statevar{maxitr} is very small,
such as zero or one.

\taossubhead{Initial Estimates and Limits on Parameters}

The method requires an initial estimate, or starting trajectory.  It works best when
the starting trajectory has characteristics similar to the final solution.  An initial
value must be provided for every parameter.  Optional upper and lower boundaries
and an optional normalizing reference value may also be supplied.

Values are identified by the same parameter numbers used in the problem file.  For
example,

\begin{lstlisting}[style=taosmanual]
par-3=6.5  lo-3=-10  hi-3=30
\end{lstlisting}

sets the initial estimate for parameter 3 to 6.5 and limits it to the interval
$[-10,30]$.  No reference is needed because the value is not extremely large or
small.  By contrast,

\begin{lstlisting}[style=taosmanual]
par-2=4000  lo-2=-10000  hi-2=100000  ref-2=10000
\end{lstlisting}

normalizes parameter 2 by 10,000.  A reference is normally useful when parameter
values exceed about 1000 or are smaller than about 0.01.  Initial estimates are
required; limits are optional but recommended; references are optional.  Placing the
values for each parameter on a separate line makes the block easier to read.

\taossubhead{Example Data Block}

The complete block consists of the objective, constraints, controls, and parameter
values:

\begin{lstlisting}[style=taosmanual]
*optimize a for range=max on segment 12, trajectory 1

  constrain vel=4000
  constrain gamgd=-75

  fref=100 maxitr=40 tol=1.0e-7

  par-1=9.9  lo-1=-15 hi-1=15
  par-2=8.5  lo-2=-15 hi-2=15
  par-3=5.0  lo-3=-15 hi-3=15
  par-4=1.0  lo-4=-15 hi-4=15
  par-5=65.0 lo-5=51.0 hi-5=300.0
\end{lstlisting}

This five-parameter problem maximizes range at the end of segment 12 in trajectory
1.  At the optimum, final velocity must be 4000 ft/s and the final geodetic
flight-path angle must be $-75^\circ$.  The procedure ends when it converges or after
40 iterations, with no automatic restarts.

\taossubhead{Convergence}

Optimization terminates when the maximum number of iterations is exceeded, when
convergence is detected, or when the method has difficulty making further progress.
If the iteration limit is reached, the last several objective values should be inspected.
If they are changing very little, the resulting trajectory is probably near optimum;
otherwise \statevar{maxitr} should be increased.

The message ``optimization procedure does not detect convergence, but it cannot
make further progress towards the minimum'' means that the method found a search
direction, but the one-dimensional search along that direction could not locate a
better point.  This usually means the method has converged as far as the numerical
precision and tolerance permit.  The message can be treated as a warning; in general,
it indicates that a good solution has been found.

\taossubhead{Troubleshooting the Optimization Loop}

Trajectory optimization is more an art than a science because methods such as
\texttt{vf02ad} are sensitive to the starting trajectory and the accuracy of the partial
derivatives.  Experience and judgement are required.  The loop definition is the
first item to inspect if the method fails to converge or produces a poor result.  When
a parameter is varied, the objective or at least one constraint should change.  A
parameter that affects neither can be removed.

The loop ends at the segment named by the objective.  At that point, every trajectory
and segment referenced by a constraint must already have been computed.  If a loop
involves several trajectories and \texttt{integ=0}, problems may occur; it is often
more practical to use \texttt{integ=1}, even though it is less efficient.

Test runs with \texttt{maxitr=0} are useful for checking the loop.  A few runs with
different parameter values quickly show whether the objective and constraints
respond correctly and whether the starting trajectory is reasonable.

\taossubhead{Troubleshooting Segment Final Conditions}

When a parameter controls the length of a segment, a time-based final condition is
preferred:

\begin{lstlisting}[style=taosmanual]
*when tseg>opta-3 goto 5
\end{lstlisting}

Time is strictly increasing, while quantities such as altitude may cross the same
value more than once.  Only one optimized time final condition should be used.  If
several final conditions compete, different conditions can end the segment on
different iterations and confuse the optimizer.  Additional endpoint requirements
should be expressed as optimization constraints.

\taossubhead{Troubleshooting the Starting Trajectory}

The method is sensitive to the initial parameter estimates, particularly when many
parameters are used.  If a loop fails or gives a poor solution, try different estimates;
even small changes sometimes restore convergence.  An optimum should not be
trusted completely because a local rather than global minimum can be found.  Rerun
an important analysis from a substantially different starting trajectory.  If the same
solution results, a global optimum has probably been located.

\taossubhead{Troubleshooting Numerical Derivatives}

Numerical derivatives are formed by very small perturbations of the trajectory.  The
integration step size in the \problemblock{*integ} block must be small enough to
resolve those perturbations.  If it is too large, integration errors produce incorrect
derivatives and the optimizer moves toward an incorrect solution.  Experimentation
is required to determine an adequate step.

If integration accuracy is not the problem, derivative truncation error may be too
large.  Decrease \statevar{dx} by an order of magnitude or more, or select central
differences with \texttt{derivs=0}.  The increment cannot be made arbitrarily small;
when roundoff becomes larger than truncation error, increase \statevar{dx} and use
central differences.

\taossubhead{Troubleshooting Surveys}

Optimization with surveys often causes problems because one starting solution may
not work for all survey values.  This is especially true when the parameters form an
angle-of-attack time history.  Such trajectories are best run individually unless the
survey value changes only slightly.  The restart and adjustment options can
sometimes help, but they should be treated as a last resort.

\taossubhead{Troubleshooting Angle-of-Attack Time Histories}

A common problem occurs when the final optimized abscissa becomes smaller than
the preceding fixed time:

\begin{lstlisting}[style=taosmanual]
*fly alpha vrs tseg
  opta-1    0
  opta-2   10
  opta-3   20
  opta-4  opta-5
\end{lstlisting}

If \texttt{opta-5} falls below 20, the time values are no longer strictly increasing and
table interpolation fails.  Its lower limit must therefore be set to a value slightly
larger than the preceding time point.
